% Main Document Settings
\documentclass[aspectratio=169]{beamer}
\usetheme{Goettingen}

\date{July 19, 2026}

% Enable polyglossia and alter default fonts to display polytonic Greek
\usepackage{polyglossia}
\usepackage{fontspec}
\usepackage[p, osf]{cochineal} % proportional and oldstyle Cochineal typeface
\setmainfont{cochineal}
\setsansfont{cochineal}
\setmonofont{cochineal}
% \setmainfont{FreeSerif}
% \setsansfont{FreeSerif}
% \setmonofont{FreeMono}

% Allow fontawesome icons, if desired
\usepackage{fontawesome}

% Allow code quotations
\usepackage{listings}

% Demonstrational packages for parallel typesetting
\usepackage{parallel}
%\usepackage{reledmac}
%\usepackage{reledpar}

% Allows ConTeXt and other logos
\usepackage{hologo}

\title{\LaTeX\ for Elegant Humanities Handouts \\ in Conferences and Classrooms }
\author{Corey Stephan, Ph.D. \break \href{https://www.coreystephan.com}{coreystephan.com}}

\begin{document}
\setdefaultlanguage{english}
\setotherlanguage[variant=ancient]{greek}
	
\begin{frame}[plain]
    \maketitle
\end{frame}

\begin{frame}
	\frametitle[About me]{About me}
I serve as Lecturer in Religious Studies at the University of Houston, course developer and instructor in the Th.D. program at St. Leo University, and instructor of Catholic theology in the core curriculum at the University of St. Thomas in Texas. \break

Broadly, I study patristics (the Church Fathers) and the theology of the Latin Middle Ages and Greek Byzantium. My research topics include Greek patristic Christology and cosmology, plus their Latin Scholastic reception. I am the author of \textit{Maximus the Confessor’s Thomistic Legacy} (Pontifical Institute of Medi\ae val Studies: Toronto, 2026). \break

I advocate software freedom by delivering talks like this one. \break

For more about my work, visit \href{https://www.coreystephan.com}{coreystephan.com}. \break

\end{frame}

\begin{frame}
    \frametitle[Everything]{\LaTeX\ beyond its normal bounds}
The world's best established software system for typesetting and document formatting, \LaTeX\ has been popular throughout mathematical and physical sciences for decades. \break

The utility of \LaTeX\ for humanities disciplines, however, tends to be overlooked. \break

\LaTeX\ can be used for typesetting in most instances in which a scholar of humanities needs to present her materials. \break

I prepared these slides in \LaTeX\ -- with less frustration than I would have had with a graphical presentation program -- as 1 file:\break

{\href{https://github.com/historical-theology/latex-talks/blob/main/tug/humanities-latex.tex}{github.com/historical-theology/latex-talks/blob/main/tug/humanities-latex.tex}}

\end{frame}

\begin{frame}
	\frametitle[Tutorial]{Today's tutorial}
In this tutorial, I will introduce attendees to the use of \LaTeX\ for a common task in humanities scholarship, namely, \break

typesetting texts with unusual line length and other organizational needs, such as primary texts with accompanying translations. \break

The `thought experiment' will be that of determining how to prepare a handout for a historical conference or class presentation with perfectly organized excerpts of a Byzantine Greek text and an original English translation in parallel. \break
\end{frame}

\begin{frame}
	\frametitle[\LaTeX]{As a humanities scholar, why choose \LaTeX ? {\scriptsize (1/2)} }
	
\begin{flushright}
	\textit{I intend this and other general explanations of \LaTeX\ \break and its usage that I provide to be for the benefit of my \break colleagues who might view this presentation online.}\break
\end{flushright}

In humanities disciplines, such as my own area of historical theology, editors expect us to prepare our materials for publications with bug-for-bug compatibility with Microsoft Word for back-and-forth markup. In practice, this expectation requires us to use word processors, such as (my favorite option) LibreOffice Writer. \break

A \textbf{word processor} presents material in a \textbf{W}hat \textbf{Y}ou \textbf{S}ee \textbf{I}s \textbf{W}hat \textbf{Y}ou \textbf{G}et (\textbf{WYSIWYG}, ``whizzy-wig'') style. While one writes, one is simultaneously thinking about and working on formatting.

\end{frame}

\begin{frame}
	\frametitle[\LaTeX]{As a humanities scholar, why choose \LaTeX ? {\scriptsize (2/2)} }
	
In a \textbf{typesetting language}, such as \LaTeX\ (or \hologo{ConTeXt} or Typst, either of which a humanities scholar might also choose for various reasons), one is only concerned about the \emph{structure} of what one is writing. \break

One declares that specific characters will be in certain categories of writing `things' (heading, paragraph, footnote, accent type, typeface, etc.), and one allows the program to compile one's writing as `code' into a tidy final document -- every time. \break

Preparing a \textit{curriculum vitae}, conference handout, or other document with unusual typesetting requirements becomes about the writing itself rather than the formatting.
	
\end{frame}

\begin{frame}
	\frametitle[Starting \LaTeX]{Learning \LaTeX\ as a scholar in the humanities}

2 community-built, libre guidebooks \break

\textit{The Not So Short Introduction to \LaTeX} \hfill
\href{https://github.com/oetiker/lshort}{github.com/oetiker/lshort} \break

\break \textit{Modern \LaTeX} \hfill
\href{https://github.com/mrkline/modern-latex}{github.com/mrkline/modern-latex} \break

Either guidebook may be used to begin to learn how to typeset prose, letters, ancient languages, and other materials that are relevant across the humanities -- without the strong front-loading of mathematics and physical sciences that one finds in other introductions to \LaTeX. \break

In other words, either option provides an efficient introduction to \LaTeX\ for a scholar in a humanities discipline.

\end{frame}


\begin{frame}[fragile]
	\frametitle[Simple Parallel Text]{Simple setting of parallel text in \LaTeX}
There are many ways to set parallel text in \LaTeX.\break

The simplest is to use the community package called \textbf{parallel}, which is included in \hologo{TeXLive}, \hologo{MiKTeX}, and other \LaTeX\ distrubtions. For full documentation, visit: \break 

\href{https://ctan.org/pkg/parallel}{ctan.org/pkg/parallel} \break

To start using \textbf{parallel} in any \LaTeX\ document, include the following in the preamble: \break

\begin{lstlisting}
\usepackage{parallel}
\end{lstlisting}
\end{frame}

\begin{frame}[fragile]
For example, newcomers, here is an abbreviation of the preamble to this Beamer file with the most important entries, including \textbf{parallel}: \break

{\footnotesize \begin{lstlisting}
\documentclass[aspectratio=169]{beamer}
\usetheme{Goettingen}
...

\usepackage{parallel}
...

\title{\LaTeX\ for Elegant Humanities Handouts 
	\\ in Conferences and Classrooms }
\author{Corey Stephan, Ph.D. \break 
	\href{https://www.coreystephan.com}{coreystephan.com}}

\begin{document}
\end{lstlisting}}
\end{frame}

\begin{frame}[fragile]
	\frametitle[Using parallel]{Using the \textit{parallel} package {\scriptsize (1/2)}}

Basic format \break

{\footnotesize \begin{lstlisting}
\begin{Parallel}{}<left adjustment>}{<right adjustment>}
		\ParallelLText{<left text row 1>}
		\ParallelRText{<right text row 1>}
			\ParallelPar
		\ParallelLText{<left text row 2>}
	[and so on]	
\end{Parallel}	
\end{lstlisting}}
\end{frame}

\begin{frame}[fragile]
	\frametitle[Using parallel]{Using the \textit{parallel} package {\scriptsize (2/2)}}

Example \break

{\footnotesize \begin{lstlisting}
\begin{Parallel}{0mm}{0mm}
	\ParallelLText{Ave Maria,}
	\ParallelRText{Hail Mary,}
	\ParallelPar
	\ParallelLText{gratia plena,}
	\ParallelRText{full of grace,}
	\ParallelPar
	\ParallelLText{\ldots}
	\ParallelRText{\ldots}
\end{Parallel}	
\end{lstlisting}}
\vfill
\begin{Parallel}{0mm}{0mm}
	\ParallelLText{Ave Maria,}
	\ParallelRText{Hail Mary,}
	\ParallelPar
	\ParallelLText{gratia plena,}
	\ParallelRText{full of grace,}
	\ParallelPar
	\ParallelLText{\ldots}
	\ParallelRText{\ldots}
\end{Parallel}	

\end{frame}

\begin{frame}[fragile]
	\frametitle[Advanced Parallel Text]{Advanced setting of parallel text in \LaTeX\  {\scriptsize (1/2)}}
	
For advanced scholarly setting of parallel text in \LaTeX, especially long-form (e.g. a critical edition with accompanying translation), one ought to learn the two interconnected and frequently updated community packages \break 

\textbf{reledpar} (for parallel typesetting specifically) and \break 

\textbf{reledmac} (for typesetting academic editions in general),\break 

which are bundled with \hologo{TeXLive} and other \LaTeX\ distributions. \break

\textbf{reledpar} and \textbf{reledmac} are purpose-built for long-form academic parallel texts. \break

These two packages are meant to replace older methods of academic parallel text preparation in \LaTeX\ that one might find in Web searches. \break
\end{frame}

\begin{frame}[fragile]
	\frametitle[Advanced Parallel Text]{Advanced setting of parallel text in \LaTeX\  {\scriptsize (2/2)}}

For full documentation, visit: \break

\href{https://ctan.org/pkg/reledpar}{ctan.org/pkg/\textbf{reledpar}} \break

\href{https://ctan.org/pkg/reledmac}{ctan.org/pkg/\textbf{reledmac}} \break
	
To start using \textbf{reledpar} and \textbf{reledmac}, remember them in the preamble: \break

\begin{lstlisting}
\usepackage{reledpar}

\usepackage{reledmac}
\end{lstlisting}
	
\end{frame}

\begin{frame}[fragile]
	\frametitle[Non-Latin Text]{Setting of non-Latin text in \LaTeX}

The most important packages for handling text outside of Latin alphabets are the community packages \textbf{babel} and \textbf{polyglossia}, which are included in \hologo{TeXLive} and other \LaTeX\ distributions. \break

Although it is possible to convert a project from one of these metapackages two the other, one ought to invest the time to choose carefully between either of these two at the start of a long-form project.
\end{frame}

\begin{frame}[fragile]
	\frametitle[Polyglossia vs. Babel]{polyglossia vs. babel {\scriptsize (1/2)}}
	
In \textit{The \LaTeX\ Companion: 3rd Edition}, Frank Mittelbach and Ulrike Fischer, both of whom are fellow speakers at TUG 2026, suggest that despite the fact that \textbf{babel} is older than \textbf{polyglossia}, \break 

\textbf{babel} \underline{generally is preferrable}: \break

 ``Unless you are in a special situation where a language or script is notably better supported by \textbf{polyglossia}, the recommendation is to use \textbf{babel} with all engines'' (entry 13.7.2, ``Polyglossia''). \break

Since \textbf{babel} technically supports the same languages as \textbf{polyglossia}, a situation in which one would need \textbf{polyglossia} over \textbf{babel} will be rare, but a scholar working with ancient languages might have such a rare situation. \break
\end{frame}

\begin{frame}[fragile]
	\frametitle[Polyglossia vs. Babel]{polyglossia vs. babel {\scriptsize (2/2)}}
	
In (old) web forum discussions, one might read that \textbf{babel} will work better with \hologo{LuaLaTeX}, whereas \textbf{polyglossia} will work better with \hologo{XeLaTeX}. \break 

In 2026, however, one can choose either \textbf{babel} or \textbf{polyglossia} for either of the two mainline \LaTeX\ engines that are built with full Unicode support, \hologo{XeLaTeX} or \hologo{LuaLaTeX}. \break

Results might vary depending on the specifics of a given project. \break

(I typically use \hologo{LuaLaTeX} to compile \textbf{polyglossia}, including for these slides.) \break
	
To make an informed choice, one should explore each package's documentation: \break

\href{https://www.ctan.org/pkg/babel}{ctan.org/pkg/\textbf{babel}} \break

\href{https://www.ctan.org/pkg/polyglossia}{ctan.org/pkg/\textbf{polyglossia}} \break
	
\end{frame}

\begin{frame}[fragile]
	\frametitle[Learning babel]{Learning babel}
	
To learn \textbf{babel}, I recommend Chapter 13: ``Localizing documents'' in \textit{The \LaTeX\ Companion: 3rd Edition} by Mittelbach and Fischer, \break

which succinctly introduces localization in \LaTeX\ via \textbf{babel} in a practical, didactic (teaching) style, that is,  \break

pages 297--342, in Mittelbach, Frank, and Ulrike Fischer. \textit{The LaTeX Companion: Third Edition - Part II.} Addison-Wesley, 2023.

\end{frame}

\begin{frame}[fragile]
	\frametitle[Learning polyglossia]{Learning polyglossia}
	
To learn \textbf{polyglossia}, I recommend the package's official documentation, that is, \break
	
\href{https://www.ctan.org/pkg/polyglossia}{ctan.org/pkg/\textbf{polyglossia}}
	
\end{frame}

\begin{frame}[fragile]
	\frametitle[Using polyglossia]{Starting a project with \textit{polyglossia}}

Here, I will model \textbf{polyglossia}, which is what I happen to be familiar with using for polytonic (ancient and Byzantine) Greek. {\footnotesize (Yet, after reading the recommendation of Mittelbach and Fischer, I intend to try \textbf{babel} for future projects.)} \break

Include the package in the preamble:
\begin{lstlisting}
\usepackage{polyglossia}
\end{lstlisting} 

\vspace{1em}

Then, \emph{after} the beginning of the document but \emph{before} actual type, declare the languages and their variants. For example, here is this .tex file's polyglossia declarations:

\begin{lstlisting}
\begin{document}
\setdefaultlanguage{english}
\setotherlanguage[variant=ancient]{greek}
\end{lstlisting}

\end{frame}

\begin{frame}[fragile]	
	
When using one of the `other' declared languages, toggle it when needed.

A language may be toggled for a long excerpt like so:\footnote{Here and in other code samples, the Greek text is spaced incorrectly due to a limitation of using the package \verb*|lstlisting}
	
{\footnotesize \begin{lstlisting}
\begin{greek}
<greek text>
...
\end{greek}
\end{lstlisting}}

{\footnotesize \begin{lstlisting}
\begin{greek}
ἐν ἀρχῇ ἦν ὁ λόγος, \\
καὶ ὁ λόγος ἦν πρὸς τὸν θεόν, \\
καὶ θεὸς ἦν ὁ λόγος. \\
\end{greek}
\end{lstlisting}}

\begin{greek}
	ἐν ἀρχῇ ἦν ὁ λόγος, \\
	καὶ ὁ λόγος ἦν πρὸς τὸν θεόν, \\
	καὶ θεὸς ἦν ὁ λόγος. \\
\end{greek}

\end{frame}

\begin{frame}[fragile]	
	
A language may be toggled for a single line like so: \break
	
{\footnotesize \begin{lstlisting}
\text<language>{<text in language>}
e.g.
\textgreek{<Greek text>}	\end{lstlisting}} \break
	
\vfill
This might be used to insert a keyword, e.g. \break 
{\textgreek{ὁ λόγος}} (produced with \verb|\textgreek{ὁ λόγος}|), \break
inside prose.
	
\end{frame}

\begin{frame}[fragile]
	\frametitle[Combining parallel and polyglossia]{Combining \textit{parallel} and \textit{polyglossia}}

The simplest way of showing multilingual text in parallel is \break

{\footnotesize \begin{lstlisting}
\begin{Parallel}{0mm}{0mm}
	\text<language>{\ParallelLText{<ancient text 1>}}
	\ParallelRText{<modern text>}
		\ParallelPar
	\text<language>{\ParallelLText{<ancient text 2>}}
	\ParallelRText{<modern text 2>}
\end{Parallel}	
\end{lstlisting}} \break

\vspace{1em}

but since this is tedious for more than a few lines, I remind everyone of \break \textbf{reledpar} and \textbf{reledmac} for longer and/or more complicated parallel text displays. \break

\end{frame}

\begin{frame}[fragile]

Here is a simple example of \textbf{parallel} and \textbf{polyglossia} working in combination to display my own translation from the opening chapter of \textit{Maximus the Confessor's Thomistic Legacy}:\break
	
{\footnotesize \begin{lstlisting}
\begin{Parallel}{0mm}{0mm}
	\textgreek{\ParallelLText{Νῦν συγχορεύεις ἀληθῶς}}
	\ParallelRText{Now, truly, you celebrate together}
	\ParallelPar
	\textgreek{\ParallelLText{σὺν τοῖς λοιποῖς θεολόγοις}}
	\ParallelRText{with the rest of the theologians}
...
\end{Parallel}	
	\end{lstlisting}} \break

For a conference or classroom handout, the result might look like the next slide:
	
\end{frame}

\begin{frame}[fragile]

Canon in Honor of Thomas Aquinas, Ode 3\footnote{The Greek text is from Raffaele Cantarella, ed., “Canone greco inedito di Giuseppe vescovo di Methone in onore di San Tommaso d’Aquino,” \emph{Archivum Fratrum Praedicatorum} 12 (1934): 145--85, at 155--58.} \\
{\small Translated by Corey Stephan, Ph.D.} \break

{\footnotesize \begin{Parallel}{0mm}{0mm}
	\textgreek{\ParallelLText{Νῦν συγχορεύεις ἀληθῶς}}
	\ParallelRText{Now, truly, you celebrate together}
		\ParallelPar
	\textgreek{\ParallelLText{σὺν τοῖς λοιποῖς θεολόγοις}}
	\ParallelRText{with the rest of the theologians}
		\ParallelPar
	\textgreek{\ParallelLText{καὶ τοῖς ἄλλοις θεολήπτοις ἀνδράσι,}}
	\ParallelRText{and the other divinely inspired men}
		\ParallelPar
	\ParallelRText{with the rest of the theologians}
	\textgreek{\ParallelLText{περὶ τὸν θρόνον Χριστοῦ}}
	\ParallelRText{around the throne of Christ,}
		\ParallelPar
	\textgreek{\ParallelLText{τὸν ἱερὸν ἐξᾴδοντες}}
	\ParallelRText{singing the sacred}
		\ParallelPar
	\textgreek{\ParallelLText{καὶ τρισάγιον ὕμνον}}
	\ParallelRText{and Thrice-Holy Hymn}
		\ParallelPar
	\textgreek{\ParallelLText{σὺν ταῖς ἁγίαις δυνάμεσιν.}}
	\ParallelRText{with the holy powers.}
		\ParallelPar
	\vspace{1em}
	\textgreek{\ParallelLText{Ἔνθα πατέρων ὁ χορὸς}}
	\ParallelRText{Where there is the chorus of Fathers}
		\ParallelPar
	\textgreek{\ParallelLText{καὶ ἱερῶν θεολόγων}}
	\ParallelRText{and holy theologians,}
		\ParallelPar
	\textgreek{\ParallelLText{ὡς ὑψίνους καὶ λεπτὸς θεόλογος}}
	\ParallelRText{as a high-minded and subtle theologian}
		\ParallelPar
	\textgreek{\ParallelLText{κατεσκήνωσας, Θωμᾶ,}}
	\ParallelRText{you do abide, Thomas,}
		\ParallelPar
	\textgreek{\ParallelLText{σὺν τούτοις αὐλιζόμενος,}}
	\ParallelRText{dwelling with them,}
		\ParallelPar
	\textgreek{\ParallelLText{θεολογῶν ἀπαύστως}}
	\ParallelRText{unceasingly theologizing,}
		\ParallelPar
	\textgreek{\ParallelLText{τριάδα, πάτερ, τὴν ἄκτιστον.}}
	\ParallelRText{O Father, about the Uncreated Trinity.}
		\ParallelPar
\end{Parallel}}	

\end{frame}

\begin{frame}
	\frametitle[Other \LaTeX\ software]{Other \LaTeX\ Tools for Scholars}

\textbf{BibLaTeX} -- comprehensive citation and bibliography management, capable of being interlinked with a variety of citation managers (notably, Zotero and JabRef). \break

\textbf{TeXstudio} and other \LaTeX\ -specific graphical writing environments \break
LaTeX can be written in any plain text editor, but a designated environment improves efficiency and helps avoid code errors. \break

Designated packages for individual languages, e.g. \textbf{Serto} for Syriac \break 
See Chapter 13 of \textit{The \LaTeX\ Companion: Third Edition} for other recommendations for language-specific package options, or run search queries in CTAN. \break

\textbf{microtype} for small adjustments to type, such as improved automatic hyphenation

\end{frame}

\begin{frame}

I thank the organizers of TUG 2026. \break

Now, I will take questions, comments, suggestions, and more.
	
\end{frame}

\end{document}
